\documentclass[conference,onecolumn]{IEEEtran}
\IEEEoverridecommandlockouts
\usepackage[utf8]{inputenc}
\usepackage[T1]{fontenc}
\usepackage[brazil]{babel}
\usepackage{cite}
\usepackage{amsmath,amssymb}
\usepackage{array}
\usepackage{tabularx}
\usepackage{booktabs}
\usepackage{url}
\usepackage{xcolor}
\usepackage{rotating}
\usepackage{tikz}
\usetikzlibrary{positioning,arrows.meta,babel,calc}
\usepackage[edges]{forest}
\usepackage[hidelinks]{hyperref}
\renewcommand{\tablename}{Tabela}
\renewcommand{\figurename}{Figura}
\newcolumntype{Y}{>{\raggedright\arraybackslash}X}
\newcolumntype{P}[1]{>{\raggedright\arraybackslash}p{#1}}
\newcommand{\sint}{\textsc{sintético}}
\newcommand{\HO}{\textsuperscript{[HO]}}
% status epistemológico
\newcommand{\stEV}{\textsc{ev}}
\newcommand{\stPD}{\textsc{pd}}
\newcommand{\stHO}{\textsc{ho}}
% operadores CTT
\newcommand{\enab}{\ensuremath{\gg}}
\newcommand{\enabinfo}{\ensuremath{[\,]\!\gg}}
\newcommand{\choice}{\ensuremath{[\,]}}
\newcommand{\interl}{\ensuremath{|||}}
\newcommand{\ordind}{\ensuremath{|\!=\!|}}
\newcommand{\deact}{\ensuremath{[>}}
\newcommand{\susp}{\ensuremath{|>}}
\newcommand{\iter}{\ensuremath{^{*}}}
\newcommand{\ev}[1]{\texttt{#1}}
\newcommand{\cena}[3]{\textbf{#1}\\\textit{#2}\\[-2pt]\rule{\linewidth}{0.3pt}\\#3}

\begin{document}

\title{Caderno de Especificação de Design do UXAgent (Anexo E1)\\[4pt]\large Metamensagem, paridade de signos, CTT, GOMS/KLM, MoLIC, rupturas, MAC e sistema de ajuda do protótipo CasaFácil}

\author{\IEEEauthorblockN{Maíra Beatriz de Almeida Lacerda}
\IEEEauthorblockA{\textit{Engenharia de Computação} \\
\textit{CEFET-MG -- Campus V}\\
Divinópolis, Brasil \\
maira@aluno.cefetmg.br}
\and
\IEEEauthorblockN{Maria Eduarda Teixeira Souza}
\IEEEauthorblockA{\textit{Engenharia de Computação} \\
\textit{CEFET-MG -- Campus V}\\
Divinópolis, Brasil \\
maria@aluno.cefetmg.br}
}

\maketitle

\begin{abstract}
Este caderno é o Anexo E1 da fase Desenvolvimento e Implementação do UXAgent~\cite{lara2026p3} e reúne, numa única versão do design, os artefatos exigidos pela Parte~I (Projeto de Interação) e pelos apêndices do DesignUXAgent~\cite{lara2026a3}: a metamensagem para a persona primária e o espaço de design com dois receptores (a pessoa e o agente); o esquema conceitual de signos, com a tabela de paridade que fixa como cada signo da tela é traduzido na observação $o_t$ entregue ao LLM; a reescrita em CTT das quatro tarefas herdadas da Fase~II; o CMN-GOMS de uma tarefa de calibração e a linha de base KLM de especialista, com a comparação de duas alternativas para a cena principal; o diagrama MoLIC, sua tradução em máquina de estados e as leituras de cada elemento; a matriz de prevenção e recuperação de rupturas (PP, PA, AL, RA, CE); oito etiquetas do MAC operacionalizadas como padrões de eventos; o sistema de ajuda sob demanda, organizado por tipo de dúvida; e o registro de \textit{rationale}. Cada decisão recebe o status de evidência (\stEV), princípio de design (\stPD) ou hipótese operacional (\stHO). Nenhum destes modelos é entregue ao agente: eles pertencem ao analista.
\end{abstract}

\begin{IEEEkeywords}
UXAgent, Engenharia semiótica, MoLIC, CTT, GOMS, KLM, Sistema de ajuda, Paridade de signos.
\end{IEEEkeywords}

% =====================================================================
\section{Escopo, base e versão do design}\label{sec:escopo}

\textit{Base herdada.} O design parte do Alinhamento Rastreável (Fase~II)~\cite{lacerda2026faseII}, a versão mais recente da modelagem: persona Lúcia, perfil e comparador, cenários CN1--CN4, tarefas C1 e C2 (calibração) e V1 e V2 (validação, congeladas desde 26/10/2026), HTAs completas, matriz L1--L6 com os mecanismos M1--M6 e o dicionário de métricas. A Tabela~\ref{tab:equiv} relaciona a nomenclatura da Fase~II com a da Modelagem (Prova~1), que tinha cinco tarefas.

\begin{table}[!h]
\caption{Correspondência entre tarefas da Modelagem e da Fase~II}
\label{tab:equiv}
\centering
\footnotesize
\begin{tabularx}{\textwidth}{@{}P{2.4cm}P{2.4cm}P{2.4cm}Y@{}}
\toprule
\textbf{Modelagem} & \textbf{Fase II / este caderno} & \textbf{Conjunto} & \textbf{Observação} \\
\midrule
T1 Luz apagou & C1 & Calibração & Mesmo estado inicial e critério de sucesso \\
T2 Viagem & C2 & Calibração & Idem \\
T3 Ventilador & V1 & Validação (congelada) & Idem; não usada para escolher alternativas nem ajustar a ajuda \\
T4 Ar-condicionado & V2 & Validação (congelada) & Idem \\
T5 Volta da viagem & --- & --- & Retirada na Fase II; a 2.ª exposição ao caminho de rotinas (L5) passa a ocorrer em V1 e V2 \\
\bottomrule
\end{tabularx}
\end{table}

\textit{Versão.} Este caderno define \texttt{ui\_version = casafacil-1.0} e \texttt{design\_version = e1-2026-11-22}. Toda alteração de rótulo, fluxo, ajuda ou observação depois do início da calibração gera nova versão e recomeça a calibração~\cite{lara2026a3}.

\textit{Três regras.} (i)~\textit{Design para a persona, nunca para o agente}: nenhuma decisão é justificada pelo desempenho do agente. (ii)~\textit{Equivalência de canal}: a observação $o_t$ não contém significado que a persona não obteria pela tela; quando a equivalência não é possível, a diferença é declarada como \stHO. (iii)~\textit{Rastreabilidade de design}: cada decisão relevante registra questão, opções, critérios, decisão e status (Seção~\ref{sec:rationale})~\cite{dix2004}.

\textit{Divergências encontradas e registradas.} A Tabela~\ref{tab:diverg} lista o que este caderno altera em relação às fases anteriores. As alterações atingem só elementos de design ou tarefas de calibração; V1 e V2 não foram alteradas.

\begin{table}[!h]
\caption{Divergências em relação à Modelagem e à Fase~II (registradas, não escondidas)}
\label{tab:diverg}
\centering
\footnotesize
\begin{tabularx}{\textwidth}{@{}P{3.0cm}P{4.2cm}YP{1.2cm}@{}}
\toprule
\textbf{Item} & \textbf{Antes} & \textbf{Agora e motivo} & \textbf{Status} \\
\midrule
Conteúdo da observação & ``identificador, papel, rótulo e estado'' de cada elemento (Fase II, Seç.~VII) & Só signos visíveis, com texto literal e ícones descritos pela forma; sem identificadores internos nem estado escrito por extenso quando a tela não o escreve (Seção~\ref{sec:signos}). Motivo: o identificador interno é vazamento~\cite{lara2026a3} & \stPD \\
``Excluir rotina'' & Sem confirmação (Fase II); com confirmação (Modelagem) & Com alerta de confirmação (AL) e faixa ``Desfazer'' em S2 (RA). Motivo: ação irreversível ao lado de controles frequentes; toda prevenção precisa de recuperação apoiada (Seção~\ref{sec:rupturas}) & \stPD \\
$A_{\min}$ de C1 & 6 (A) e 4 (B) & 5 (A) e 5 (B), recontados pelo CTT e pelo KLM: Concluir está na barra fixa, e a causa só é lida em S5 (A) ou no painel (B). C1 é de calibração, portanto pode ser corrigida & \stPD \\
KLM & Só contagem de operadores, sem tempos & Contagem e tempo de especialista com tabela única de tempos (Seção~\ref{sec:klm}); o tempo é linha de base do analista, e não tempo comportamental da persona & \stHO \\
Esquema de eventos & 17 tipos & Acrescenta \ev{dialog\_shown} e \ev{dialog\_choice} (AL e CE), campos de camada e cena de origem em \ev{support\_open}, e os eventos derivados \ev{method\_selected} e \ev{communicability\_tagged} (Seção~\ref{sec:eventos}) & \stPD \\
Ajuda & Tela S6 única & Duas camadas sob demanda (dica local ``?'' em S3 e S4; tela S6) mais o painel ``Por quê?'' em B (Seção~\ref{sec:ajuda}) & \stPD \\
Estado dos aparelhos & Interruptor (só gráfico) & Interruptor acompanhado do texto ``Ligada''/``Desligada'' sob o nome do aparelho, para não depender só de cor e posição & \stPD \\
\bottomrule
\end{tabularx}
\end{table}

% =====================================================================
\section{Metamensagem e espaço de design com dois receptores}\label{sec:meta}

\subsection{Metamensagem para a persona primária}
Redigida no modelo de de Souza~\cite{desouza2005}, \cite{barbosa2010}, \cite{aula14}. A primeira parte vem do perfil e da persona da Fase~II; a segunda, deste design.

\begin{quote}\small
\textit{Este é o meu entendimento, como designer, de quem você é:} uma pessoa de 60 anos ou mais, que usa o celular todos os dias para conversar, ver fotos e fazer chamadas de vídeo, mas que não criou as automações da casa --- quem as criou foi outra pessoa, como sua filha. Você lê com atenção em português, prefere resolver as coisas sozinha e tem receio de apagar algo por engano. \textit{Do que você quer ou precisa fazer:} entender por que um aparelho mudou sem que você tocasse nele e conseguir desfazer ou impedir essa mudança, sem estragar o que foi programado. \textit{De como prefere fazê-lo e por quê:} pelo celular, com poucos toques e sem digitar, conferindo cada resultado na tela e recorrendo à Ajuda quando não reconhece um símbolo, porque assim continua no controle da própria casa. \textit{Este, portanto, é o sistema que projetei para você, e esta é a forma como você pode ou deve utilizá-lo:} na tela Início você vê seus aparelhos, cada um com o estado escrito, e um aviso sempre que algo muda sozinho; no Histórico você vê quem fez cada mudança --- você ou uma rotina --- e, na versão~B, pode tocar em ``Por quê?'' no próprio aviso; em Rotinas você pode pausar, editar ou excluir as regras automáticas, e nenhuma exclusão acontece sem a sua confirmação, podendo ainda ser desfeita; a Ajuda, em toda tela, explica cada símbolo e cada passo, mas só aparece quando você pede.
\end{quote}

\textit{Restrição declarada.} Ícones sem texto (pausar e editar) e o rótulo em inglês ``Smart Schedule'' contrariam a própria metamensagem, mas são mantidos de propósito: reproduzem aplicativos comerciais~\cite[pp.~34--35]{tsuchiya2018} e são a condição que exercita L2. A metamensagem compensa com a promessa de ajuda sob demanda, e não pela troca dos signos (decisão D12, Seção~\ref{sec:rationale}).

\subsection{Espaço de design com dois receptores}
A Tabela~\ref{tab:jakobson} aplica os elementos do modelo de Jakobson adotados pela engenharia semiótica~\cite{barbosa2010}. É nela que se decide a equivalência de canal.

\begin{table}[!h]
\caption{Espaço de design com dois receptores}
\label{tab:jakobson}
\centering
\scriptsize
\begin{tabularx}{\textwidth}{@{}P{1.2cm}P{2.6cm}YYP{3.6cm}@{}}
\toprule
\textbf{Elemento} & \textbf{Pergunta do designer} & \textbf{Para a persona} & \textbf{Para o agente} & \textbf{Diferença declarada e status} \\
\midrule
Emissor & Que limitações e intenções do designer devem ser comunicadas? & O aplicativo mostra o que a casa faz sozinha e permite desfazer; as rotinas foram criadas por outra pessoa; exclusão exige confirmação & As mesmas intenções, reconstruídas só a partir dos signos de $o_t$; nenhuma explicação adicional do designer no prompt & Nenhuma no conteúdo. O agente recebe também a instrução de formato da resposta (escolher um número da lista), que não comunica intenção do designer da interface. \stPD \\
\addlinespace
Receptor & Que aspectos do usuário, como interpretado pelo designer, devem ser comunicados? & Novata em automação, lê português, receio de errar, usa a Ajuda quando insegura & LLM com a ficha da persona e o estado $x_t$ (crença, notas, rótulos aprendidos); comportamento do perfil vem de M1--M6, não do texto & O LLM conhece convenções de aplicativos pelo pré-treinamento (por exemplo, o ícone de duas barras). M2 compensa em parte, com probabilidade $q_{\mathrm{rot}}$. \stHO \\
\addlinespace
Contexto & Que elementos do contexto o sistema processa, e como? & Casa, à tarde, uma mão, sozinha; o aplicativo processa só os estados dos aparelhos e das rotinas & Sem contexto físico; o contexto chega apenas pela instrução literal da tarefa & Sem tempo, cansaço, pressão nem interrupções; diferença declarada nos limites de uso. \stHO \\
\addlinespace
Código & Qual é a linguagem de interface? & Manipulação direta por toque; texto em português; dois ícones sem texto; um rótulo em inglês & Texto em português que descreve a tela na ordem visual; ícones descritos pela forma; ações numeradas & Tradução visual $\rightarrow$ verbal. Descrever a forma do ícone é tomado como equivalente a vê-lo. \stHO \\
\addlinespace
Canal & Por quais canais a metacomunicação chega? & Tela do smartphone (signos estáticos e dinâmicos); Ajuda, dica local e painel ``Por quê?'' (metalinguísticos), só quando abertos & $o_t$ textual a cada passo; os signos metalinguísticos só entram em $o_t$ depois de uma ação registrada de abrir ajuda & Canal diferente, conteúdo equivalente (regra de equivalência de canal). \stPD \\
\addlinespace
Mensagem & O que se quer contar, e com que efeito? & A metamensagem acima: ``você pode ver e desfazer o que a casa faz sozinha'' & A mesma, inferida pelo agente a partir dos signos & Nenhuma intencional; se o agente a reconstruir melhor que a persona, o efeito vem do LLM, não do design. \stHO \\
\bottomrule
\end{tabularx}

\end{table}

% =====================================================================
\section{Esquema conceitual de signos e paridade}\label{sec:signos}

\subsection{Telas do protótipo}
A Tabela~\ref{tab:telas} lista os elementos de cada tela na ordem visual (de cima para baixo, da esquerda para a direita). Essa ordem é a mesma da observação $o_t$, e a vizinhança entre elementos é a usada por M3 (deslize), calculada pelo leiaute e não pelo texto.

\begin{table}[!h]
\caption{Inventário de telas (ordem visual = ordem em $o_t$; \textit{[B]}: só na versão B)}
\label{tab:telas}
\centering
\scriptsize
\begin{tabularx}{\textwidth}{@{}P{2.4cm}Y@{}}
\toprule
\textbf{Tela} & \textbf{Elementos, em ordem} \\
\midrule
Barra superior (todas) & Botão ``Concluir''; botão ``Desistir'' \\
S1 Início & Aviso (se houver mudança automática): ``<Aparelho>: <ligada|desligada> às <hh:mm>'' \textit{[B: + botões ``Por quê?'' e ``Desfazer'']}; título do cômodo; para cada aparelho: nome, texto do estado (``Ligada''/``Desligada''), interruptor. Aparelhos: Luz da sala, Ar-condicionado do quarto, Cafeteira, Tomada do ventilador \\
S2 Rotinas & Título ``Rotinas''; para cada rotina: nome e situação (``Ativa''/``Pausada''); faixa ``Rotina <nome> excluída. Desfazer'' (só logo após uma exclusão, até sair de S2) \\
S3 Detalhe da rotina & Botão ``Voltar''; nome; situação; regra ``Se <condição>, então <ações>''; botão com ícone de duas barras verticais (pausar) ou de triângulo (retomar), sem texto; botão com ícone de lápis (editar), sem texto, adjacente; botão ``?''; botão de texto vermelho ``Excluir rotina'' \\
S4 Editar rotina & Botão ``Cancelar''; ``Smart Schedule'' com o horário (``07h00''), botões ``$-$'' e ``$+$'' adjacentes e a nota ``de 00h a 23h''; ``Aparelhos'': caixa de seleção e nome de cada aparelho; botão ``?''; botão ``Salvar'' \\
S5 Histórico & Título; linhas ``<hh:mm> · <aparelho> · <novo estado> · <origem>'', com origem ``Você'' ou ``Rotina <nome>'', da mais recente à mais antiga \\
S6 Ajuda & Lista de perguntas (Tabela~\ref{tab:ajuda}); ao tocar, o texto da resposta; botão ``Voltar'' \\
S7 Concluir & Pergunta ``Por que isso aconteceu?'' com as opções ``Eu ou outra pessoa'', ``Defeito'', uma por rotina existente, ``Não sei'' (só em C1, V1 e V2); em C2, o texto ``Você terminou?''; botões ``Confirmar'' e ``Voltar'' \\
S8 Painel ``Por quê?'' \textit{[B]} & Texto da causa (``A rotina <nome> <ação> <aparelho> às <hh:mm> porque <condição>''); botões ``Desfazer'', ``Ir para a rotina'' e ``Fechar'' \\
S9 Dica local & Balão sobre S3 ou S4 com o significado dos ícones e campos da tela; botão ``Fechar'' \\
Diálogos & AL: ``Excluir a rotina <nome>? Ela sai da lista de rotinas.'' [Cancelar] [Excluir]. CE1: ``Você mudou esta rotina. Salvar as mudanças?'' [Salvar] [Sair sem salvar]. CE2 (em S7): ``Escolha uma resposta antes de confirmar.'' \\
\bottomrule
\end{tabularx}

\end{table}

\textit{Sem elementos desabilitados.} O CasaFácil não tem botões esmaecidos: todo elemento visível é acionável. Um botão desabilitado obrigaria a escolher entre escrever ``desabilitado'' em $o_t$ (vazamento) ou omiti-lo da lista de ações (o agente saberia que ele não funciona sem tentar). Onde seria natural desabilitar (Confirmar sem resposta em S7), usa-se captura de erro (CE2). Decisão D2.

\subsection{Regras de tradução da tela para $o_t$}
\begin{enumerate}
\item \textbf{Texto literal.} Todo texto visível entra entre aspas, como está na tela, inclusive ``Smart Schedule'' em inglês.
\item \textbf{Ícones pela forma.} Ícone sem texto vira a descrição de sua forma (``ícone de duas barras verticais, sem texto''), nunca sua função (``pausar''). Se M2 dispara, a descrição vira ``[ícone sem nome]'' e o elemento sai da lista de ações até a leitura da ajuda correspondente (L2).
\item \textbf{Só a tela atual.} Nada de outras telas, do estado JSON, de nomes de variáveis (\texttt{luz\_sala}), da origem de uma mudança quando a tela não a mostra, nem de elementos ocultos ou fora da tela.
\item \textbf{Signos dinâmicos como mudança.} Quando um signo muda entre dois passos na mesma tela, $o_t$ traz a linha ``Mudou na tela: <signo>: <antes> $\rightarrow$ <depois>'', porque a pessoa vê a mudança acontecer.
\item \textbf{Metalinguísticos sob demanda.} Textos de Ajuda, dica local e painel ``Por quê?'' só entram em $o_t$ depois da ação de abri-los, que gera \ev{support\_open}.
\item \textbf{Ações numeradas.} A lista de ações traz um número por elemento acionável na ordem visual, com verbo genérico (``tocar'', ``alternar'', ``marcar'') e a mesma descrição do elemento. Os números valem só para aquela observação; não são identificadores internos.
\item \textbf{Sem gabarito.} $o_t$ e o prompt nunca contêm HTA, CTT, métodos GOMS, $A_{\min}$, critérios de sucesso, a solução, a lista de tarefas ou dados de validação.
\end{enumerate}

\subsection{Esquema conceitual de signos estendido (tabela de paridade)}
As Tabelas~\ref{tab:signos1} e~\ref{tab:signos2} preenchem o esquema de Barbosa e Silva~\cite{barbosa2010}, estendido com as colunas \textit{Na observação do agente} (regra de equivalência de canal) e \textit{Evento} (ligação com o registro). Origem: \textit{dom}ínio, \textit{transf}ormado, \textit{aplic}ação ou \textit{conv}encional. Tipo: \textit{E}stático, \textit{D}inâmico ou \textit{M}etalinguístico.

\begin{table}[!h]
\caption{Esquema de signos e paridade --- S1, S5, S7, S8 e barras}
\label{tab:signos1}
\centering
\scriptsize
\setlength{\tabcolsep}{3pt}
\begin{tabularx}{\textwidth}{@{}P{2.2cm}P{0.9cm}P{2.6cm}P{2.2cm}P{0.5cm}YP{2.3cm}@{}}
\toprule
\textbf{Signo} & \textbf{Origem} & \textbf{Conteúdo (tipo, restrições, default)} & \textbf{Expressão} & \textbf{Tipo} & \textbf{Na observação do agente} & \textbf{Evento} \\
\midrule
aparelho.nome & dom & texto; 4 aparelhos fixos & rótulo à esquerda & E & \texttt{"Luz da sala"} & \ev{screen\_view} \\
aparelho.estado & dom & \{Ligada, Desligada\}; inicial pela tarefa & texto sob o nome + interruptor & D & \texttt{"Desligada"}, interruptor; mudança: ``Mudou na tela: Luz da sala: Desligada $\rightarrow$ Ligada'' & \ev{state\_changed} \\
aviso.mudança & aplic & ``<aparelho>: <estado> às <hh:mm>''; sem origem & faixa no topo de S1 & D & texto literal; nunca a origem & \ev{screen\_view} \\
aviso.porque \textit{[B]} & aplic & botão & texto ``Por quê?'' & E & ``tocar \texttt{"Por quê?"}'' & \ev{support\_open} (por\_que) \\
aviso.desfazer \textit{[B]} & aplic & botão; desfaz a última mudança automática & texto ``Desfazer'' & E & ``tocar \texttt{"Desfazer"}'' & \ev{state\_changed} (origem=pessoa) \\
painel.causa \textit{[B]} & aplic & texto; nome da rotina, ação, hora, condição & painel sob o aviso & M & texto literal, só após tocar ``Por quê?'' & \ev{support\_open}/\ev{close} \\
historico.linha & transf & hora, aparelho, estado, origem \{Você, Rotina <nome>\} & lista em S5 & D & linhas literais, só com S5 aberta & \ev{screen\_view}(S5) \\
barra.navegação & conv & Início, Rotinas, Histórico, Ajuda & abas com ícone e texto & E & ``Barra inferior: \texttt{"Início"}, \texttt{"Rotinas"}, \texttt{"Histórico"}, \texttt{"Ajuda"}'' (o texto basta; ícones das abas omitidos por redundância) & \ev{screen\_view} \\
barra.concluir, barra.desistir & conv (instrumento) & botões; sempre visíveis & texto & E & literal & \ev{screen\_view}(S7); \ev{task\_end} \\
concluir.pergunta & aplic & opções fechadas; nenhuma pré-marcada & botões de opção & E & ``\texttt{"Por que isso aconteceu?"}'' e opções literais; marcada: ``(selecionada)'' porque a tela mostra o círculo preenchido & \ev{cause\_selected} \\
concluir.CE2 & aplic & mensagem de captura & texto vermelho sob a pergunta & M & literal, só depois de Confirmar sem resposta & \ev{dialog\_shown}(CE) \\
\bottomrule
\end{tabularx}

\end{table}

\begin{table}[!h]
\caption{Esquema de signos e paridade --- S2, S3, S4, ajuda e diálogos}
\label{tab:signos2}
\centering
\scriptsize
\setlength{\tabcolsep}{3pt}
\begin{tabularx}{\textwidth}{@{}P{2.2cm}P{0.9cm}P{2.6cm}P{2.2cm}P{0.5cm}YP{2.3cm}@{}}
\toprule
\textbf{Signo} & \textbf{Origem} & \textbf{Conteúdo (tipo, restrições, default)} & \textbf{Expressão} & \textbf{Tipo} & \textbf{Na observação do agente} & \textbf{Evento} \\
\midrule
rotina.nome & dom (criado pela filha) & texto & item de lista; título em S3 & E & literal & \ev{screen\_view} \\
rotina.situação & aplic & \{Ativa, Pausada\} & texto sob o nome & D & literal; mudança como ``Mudou na tela'' & \ev{state\_changed}(status) \\
rotina.regra & transf & ``Se <condição>, então <ações>'' & texto em S3 & E & literal & \ev{screen\_view}(S3) \\
rotina.pausar/retomar & aplic & alterna a situação & ícone (duas barras $\leftrightarrow$ triângulo), sem texto & E/D & ``botão com ícone de duas barras verticais, sem texto''; com M2: ``[ícone sem nome]'', fora da lista & \ev{action\_executed}; \ev{state\_changed} \\
rotina.editar & aplic & abre S4 & ícone de lápis, sem texto, adjacente ao anterior & E & ``botão com ícone de lápis, sem texto''; M2 como acima & \ev{screen\_view}(S4) \\
rotina.excluir & aplic & ação destrutiva com AL e RA & texto vermelho ``Excluir rotina'' & E & literal (a cor não é descrita: o texto já diz a ação) & \ev{dialog\_shown}(AL) \\
alerta.exclusão (AL) & aplic & confirmação & diálogo modal & M & texto e botões literais & \ev{dialog\_shown}, \ev{dialog\_choice} \\
faixa.desfazer & aplic & RA da exclusão; some ao sair de S2 & faixa em S2 & D & literal enquanto visível & \ev{state\_changed}(restaurada) \\
rotina.horário & aplic & 00h--23h, passo 1~h, sem volta de 23h para 0h (PA); default da rotina & ``Smart Schedule'' + valor + ``$-$''/``$+$'' + ``de 00h a 23h'' & D & literal, em inglês; valor atualizado como ``Mudou na tela'' & \ev{action\_executed} \\
rotina.aparelhos & dom & caixas; marcadas conforme a rotina & caixa + nome & D & ``caixa com marca de seleção, \texttt{"Cafeteira"}'' ou ``caixa vazia, \texttt{"Cafeteira"}'' (a marca é visível) & \ev{action\_executed} \\
editar.salvar/cancelar & conv & botões & texto & E & literal & \ev{state\_changed}; \ev{dialog\_shown}(CE) \\
captura.CE1 & aplic & alterações não salvas & diálogo modal & M & literal & \ev{dialog\_shown}, \ev{dialog\_choice} \\
dica local ``?'' & aplic & abre S9 sobre a tela & botão ``?'' & M & ``tocar \texttt{"?"}''; o balão só após o toque & \ev{support\_open}(dica) \\
ajuda.tópico & aplic & perguntas e respostas (Tabela~\ref{tab:ajuda}) & lista em S6 & M & literal, só após abrir S6 e o tópico & \ev{support\_open}(ajuda, tópico) \\
caderno.nota & --- (artefato da persona) & caminho anotado & papel, fora da tela & --- & ``Sua anotação: ...'' com prob. $r_{\mathrm{nota}}$ (M5) & \ev{note\_shown} (só sintético) \\
\bottomrule
\end{tabularx}

\end{table}

\textit{Signos de aplicação} (rotina, situação, Smart Schedule, ícones, ``Por quê?'') só fazem sentido dentro do sistema e por isso têm explicação completa na ajuda (Seção~\ref{sec:ajuda})~\cite{barbosa2010}.

\subsection{Exemplo de observação}
Observação de S3, tarefa C2 (calibração), versão A, M2 inativo neste passo:
\begin{quote}\footnotesize\ttfamily
Tela: "Bom dia"\\
Barra superior: "Concluir", "Desistir"\\
Situação: "Ativa"\\
Regra: "Se for 07h00, então ligar Cafeteira e Luz da sala"\\
Barra inferior: "Início", "Rotinas", "Histórico", "Ajuda"\\
Ações disponíveis:\\
\ \ [1] tocar "Concluir" \ [2] tocar "Desistir" \ [3] tocar "Voltar"\\
\ \ [4] tocar botão com ícone de duas barras verticais, sem texto\\
\ \ [5] tocar botão com ícone de lápis, sem texto \ [6] tocar "?"\\
\ \ [7] tocar "Excluir rotina" \ [8] tocar "Início" \ [9] tocar "Rotinas"\\
\ \ [10] tocar "Histórico" \ [11] tocar "Ajuda"
\end{quote}
Resposta exigida do LLM: \texttt{\{"acao": 4, "justificativa": "..."\}}. O controlador valida o número contra a lista; fora da lista, pede uma nova resposta uma vez e, se falhar, registra \texttt{F-FMT}. A justificativa é traço qualitativo e não entra em métrica.

\subsection{Fronteira do contexto do LLM}
\begin{table}[!h]
\caption{O que entra e o que não entra no contexto do LLM}
\label{tab:fronteira}
\centering
\footnotesize
\begin{tabularx}{\textwidth}{@{}YY@{}}
\toprule
\textbf{Entra (a cada chamada)} & \textbf{Nunca entra} \\
\midrule
Instrução fixa de formato (escolher um número da lista e justificar em uma frase); ficha da persona (idade, experiência digital, objetivos de vida e de experiência); instrução literal da tarefa, igual à lida para pessoas; $o'_t$ (observação após M1, M2 e M5); de $x_t$: crença atual sobre a causa (se houver), notas do caderno mostradas por M5, rótulos já aprendidos e as últimas oito ações que a própria persona executou na tarefa, descritas pelo texto visível do elemento tocado (memória de trabalho, \stHO) & HTA, CTT, métodos e regras de seleção GOMS, KLM e $A_{\min}$; critério de sucesso e estado-alvo; estado JSON da interface e nomes de variáveis; outras telas; histórico da conversa (a continuidade existe só em $x_t$); limite de passos; outras tarefas, em especial V1 e V2; resultados de execuções anteriores; exemplos de execução; qualquer instrução de encenação (``aja como idosa'', ``hesite'', ``erre'') \\
\bottomrule
\end{tabularx}
\end{table}

\textit{Diferenças residuais (\stHO).} (a) O agente recebe o texto inteiro da tela a cada passo; a pessoa pode não ler tudo. (b) A forma descrita de um ícone não reproduz a percepção visual. (c) O caderno é físico para pessoas e não é registrado; para o agente é mecanismo registrado (M5). (d) O LLM pode trazer do pré-treinamento convenções de aplicativos que a persona não tem.

% =====================================================================
\section{Modelos de tarefa: da HTA ao CTT e ao GOMS}\label{sec:tarefas}

A HTA da Fase~II~\cite{annett1967} diz a lógica da tarefa; o CTT acrescenta a ordem temporal e a divisão entre persona e sistema~\cite{paterno2000}, \cite[seç.~18.3]{dix2004}; o GOMS acrescenta o conhecimento procedimental e a previsão de tempo de um usuário experiente~\cite{card1983}, \cite[cap.~12]{dix2004}, \cite{aula21}.

\subsection{CTT}
\textit{Notação.} Tipos: [A] abstrata, [U] do usuário (operação interna, não observável na Fase~II), [I] interativa (gera evento de interação), [S] do sistema (gera evento do sistema). Operadores: \enab{} habilitação; \enabinfo{} habilitação com passagem de informação; \choice{} escolha; \ordind{} independência de ordem; \interl{} intercalação; \deact{} desativação; \susp{} suspensão e retomada; \iter{} iteração; $[T]$ tarefa opcional. Pela prioridade dos operadores (\choice{} $>$ \ordind{} $>$ \interl{} $>$ \deact{} $>$ \susp{} $>$ \enab)~\cite{paterno2000}, parênteses só aparecem onde a leitura mudaria.

\textit{Convenções comuns às quatro tarefas.} A consulta à ajuda é suspensão e retomada: \textit{Núcleo} \susp{} \textit{ConsultarAjuda}, com \textit{ConsultarAjuda} = TocarAjudaOu?[I] \enab{} ExibirAjuda[S] \enabinfo{} LerAjuda[U] \enab{} Voltar[I]. A desistência e o limite do protocolo são desativação: (\textit{Núcleo} \susp{} \textit{ConsultarAjuda}) \deact{} \textit{Encerrar}, com \textit{Encerrar} = TocarDesistir[I] \choice{} AtingirLimite[S]. O CTT não representa prevenção e tratamento de erros~\cite{barbosa2010}; eles estão na Seção~\ref{sec:rupturas}.

\begin{figure}[!h]
\centering
\resizebox{\textwidth}{!}{%
\begin{forest}
for tree={grow'=east, parent anchor=east, child anchor=west, anchor=west,
  draw, rounded corners=2pt, font=\scriptsize, align=left, inner sep=2pt,
  l sep=4mm, s sep=0.6mm, edge={thin}},
op/.style={draw=none, no edge, font=\footnotesize, inner sep=0pt},
abs/.style={fill=gray!20},
usr/.style={fill=yellow!25},
int/.style={fill=blue!12},
sys/.style={fill=green!15},
[{C1 Restaurar a luz\\e explicar [A]}, abs
  [{Executar C1 [A]}, abs
    [{Núcleo C1 [A]}, abs
      [{Identificar causa [A]}, abs
        [{Pelo Histórico [A]}, abs
          [{Tocar Histórico [I]}, int]
          [\enab, op]
          [{Exibir Histórico [S]}, sys]
          [\enabinfo, op]
          [{Ler origem [U]}, usr]
        ]
        [\choice, op]
        [{Pelo painel \textit{[B]} [A]}, abs
          [{Tocar ``Por quê?'' [I]}, int]
          [\enab, op]
          [{Exibir painel [S]}, sys]
          [\enabinfo, op]
          [{Ler explicação [U]}, usr]
        ]
        [\enabinfo, op]
        [{Decidir causa [U]}, usr]
      ]
      [\ordind, op]
      [{Acender luz [A]}, abs
        [{Pelo interruptor [A]}, abs
          [{[Ir para Início] [I]}, int]
          [\enab, op]
          [{Tocar interruptor [I]}, int]
          [\enab, op]
          [{Atualizar estado [S]}, sys]
        ]
        [\choice, op]
        [{Por Desfazer \textit{[B]} [A]}, abs
          [{Tocar Desfazer [I]}, int]
          [\enab, op]
          [{Atualizar estado [S]}, sys]
        ]
      ]
      [\enabinfo, op]
      [{Concluir [A]}, abs
        [{Tocar Concluir [I]}, int]
        [\enab, op]
        [{Exibir pergunta [S]}, sys]
        [\enab, op]
        [{Escolher causa [I]}, int]
        [\enab, op]
        [{Tocar Confirmar [I]}, int]
        [\enab, op]
        [{Registrar fim [S]}, sys]
      ]
    ]
    [\susp, op]
    [{Consultar ajuda [A]}, abs
      [{Tocar Ajuda ou ``?'' [I]}, int]
      [\enab, op]
      [{Exibir ajuda [S]}, sys]
      [\enabinfo, op]
      [{Ler ajuda [U]}, usr]
      [\enab, op]
      [{Voltar [I]}, int]
    ]
  ]
  [\deact, op]
  [{Encerrar [A]}, abs
    [{Tocar Desistir [I]}, int]
    [\choice, op]
    [{Atingir limite [S]}, sys]
  ]
]
\end{forest}}

\caption{Árvore CTT de C1 (calibração), versões A e B. Cinza: abstrata; amarelo: usuário; azul: interativa; verde: sistema. Os operadores ficam entre irmãos.}
\label{fig:ctt-c1}
\end{figure}

\begin{table}[!h]
\caption{CTT de C2, V1 e V2 em notação textual (V1 e V2 reescrevem a HTA congelada sem alterá-la)}
\label{tab:ctt}
\centering
\scriptsize
\begin{tabularx}{\textwidth}{@{}P{1.1cm}Y@{}}
\toprule
\textbf{Tarefa} & \textbf{Expressão CTT} \\
\midrule
C2 & C2 = (Núcleo \susp{} ConsultarAjuda) \deact{} Encerrar \\
& Núcleo = EncontrarRotina \enab{} SuspenderAção \enab{} ConcluirC2 \\
& EncontrarRotina = TocarRotinas[I] \enab{} ExibirLista[S] \enabinfo{} (IdentificarRotina[U] \enabinfo{} TocarRotina[I] \enab{} ExibirDetalhe[S])\iter \quad \textit{(ciclo do plano 1 da HTA)} \\
& SuspenderAção = Pausar \choice{} EditarDesmarcar \\
& Pausar = TocarÍconePausa[I] \enab{} AtualizarSituação[S] \\
& EditarDesmarcar = TocarÍconeEditar[I] \enab{} ExibirEdição[S] \enab{} DesmarcarCafeteira[I] \enab{} TocarSalvar[I] \enab{} SalvarRotina[S] \\
& ConcluirC2 = TocarConcluir[I] \enab{} ExibirConfirmação[S] \enab{} TocarConfirmar[I] \enab{} RegistrarFim[S] \\
\midrule
V1 & V1 = (Núcleo \susp{} ConsultarAjuda) \deact{} Encerrar \\
& Núcleo = IdentificarCausa \enabinfo{} LocalizarRotina \enab{} ImpedirRepetição \enab{} Concluir \quad \textit{(IdentificarCausa e Concluir como em C1)} \\
& LocalizarRotina = TocarIrParaRotina[I]\textit{[B]} \choice{} (TocarRotinas[I] \enab{} ExibirLista[S] \enabinfo{} TocarHorárioDePico[I]) \\
& ImpedirRepetição = Pausar \choice{} EditarDesmarcar(Tomada do ventilador) \\
\midrule
V2 & V2 = (Núcleo \susp{} ConsultarAjuda) \deact{} Encerrar \\
& Núcleo = IdentificarCausa \enabinfo{} AbrirEdição \enab{} AjustarHorário \enab{} TocarSalvar[I] \enab{} SalvarRotina[S] \enab{} Concluir \\
& AbrirEdição = (TocarIrParaRotina[I]\textit{[B]} \choice{} (TocarRotinas[I] \enab{} TocarClimatização[I])) \enab{} TocarÍconeEditar[I] \enab{} ExibirEdição[S] \\
& AjustarHorário = (TocarMais[I] \enab{} ExibirHorário[S] \enabinfo{} ConferirHorário[U])\iter \quad \textit{(até ler 21h00; ciclo do plano 3 da HTA)} \\
\bottomrule
\end{tabularx}

\end{table}

\textit{Divergências HTA $\rightarrow$ CTT (registradas, não corrigidas).} (1)~O plano 0 de C1 é chamado de ``paralelo/ordem livre'' na Fase~II; o CTT usa \ordind{} (ordem livre sem intercalação), mas o protótipo também admite intercalar (abrir o Histórico, voltar e acender, voltar ao Histórico), e o verificador aceita essa sequência, pois julga só o estado final. (2)~Em V1 e V2, \textit{IdentificarCausa} habilita a localização com passagem de informação (\enabinfo), porque o nome da rotina lido no Histórico ou no painel é o que permite escolhê-la; a HTA congelada diz apenas ``sequência fixa''. (3)~Na HTA, ``Abrir Ajuda'' aparece como opcional dentro do subobjetivo 1 de C1; no CTT a consulta é global (\susp), pois a Ajuda está em todas as telas. Nenhuma divergência altera objetivos, operações ou critérios de sucesso de V1 e V2.

\subsection{CMN-GOMS de C2 (calibração)}
Notação de Card, Moran e Newell~\cite{card1983}, \cite[cap.~12]{dix2004}; operadores em maiúsculas.

\begin{quote}\footnotesize\ttfamily
GOAL: IMPEDIR-QUE-A-CAFETEIRA-LIGUE-AMANHÃ\\
.\ GOAL: ABRIR-LISTA-DE-ROTINAS\\
.\ .\ TOCAR "Rotinas"; VERIFICAR lista exibida\\
.\ GOAL: ABRIR-ROTINA-DA-CAFETEIRA\\
.\ .\ LOCALIZAR rotina cuja regra inclui a Cafeteira ("Bom dia")\\
.\ .\ TOCAR nome da rotina; VERIFICAR regra exibida\\
.\ GOAL: SUSPENDER-A-AÇÃO\\
.\ .\ [select\ METHOD: PAUSAR-ROTINA\\
.\ .\ .\ \ \ \ \ \ TOCAR ícone de duas barras; VERIFICAR situação "Pausada"\\
.\ .\ .\ \ \ \ METHOD: EDITAR-E-DESMARCAR\\
.\ .\ .\ \ \ \ \ \ TOCAR ícone de lápis; VERIFICAR tela de edição\\
.\ .\ .\ \ \ \ \ \ TOCAR caixa "Cafeteira"; VERIFICAR caixa vazia\\
.\ .\ .\ \ \ \ \ \ TOCAR "Salvar"; VERIFICAR regra sem a Cafeteira]\\
.\ GOAL: ENCERRAR-TAREFA\\
.\ .\ TOCAR "Concluir"; TOCAR "Confirmar"
\end{quote}
\textit{Regras de seleção para SUSPENDER-A-AÇÃO.} R1: se a pessoa quer que os demais aparelhos da rotina continuem agindo, use EDITAR-E-DESMARCAR. R2: caso contrário, use PAUSAR-ROTINA (menos operadores). \textit{Escolha observável.} O método fica visível no registro: \ev{state\_changed(status=Pausada)} indica PAUSAR; \ev{state\_changed(aparelhos)} indica EDITAR. O verificador grava o evento derivado \ev{method\_selected} no fim da tarefa; ele vira indicador, e não métrica de sucesso. Regras análogas valem para C1 (se o aviso tem ``Por quê?'', usar o painel; senão, o Histórico) e geram \ev{method\_selected} \{historico, painel\} e \{interruptor, desfazer\}.

\subsection{Linha de base KLM de especialista}\label{sec:klm}
\textit{Tabela única de tempos}~\cite{card1980}: P = 1,10~s (apontar); K = 0,28~s (pressionar; digitador médio); M = 1,35~s (preparação mental); H não ocorre (uso com uma mão, dedo já sobre a tela); R = 0 (protótipo local; a latência do LLM é desempenho do sistema e não entra em R). Um toque é modelado como P + K = 1,38~s; essa adaptação do KLM de teclado e mouse para toque é \stHO.

\textit{Regras de colocação de M, iguais para todas as alternativas}~\cite{card1980}, \cite{preece2002}: (M-a) uma M antes do primeiro toque após a exibição de tela ou painel novo (ler e decidir); (M-b) uma M antes do toque que inicia um subobjetivo da HTA, se M-a não tiver colocado uma M ali; (M-c) toques repetidos no mesmo botão formam um bloco, com uma M antes do primeiro e uma M ao final para conferir o valor; (M-d) não há M entre escolher a opção e Confirmar na mesma tela.

\begin{table}[!h]
\caption{Linha de base KLM de especialista por tarefa e versão (caminho mínimo, sem erros)}
\label{tab:klm}
\centering
\scriptsize
\begin{tabularx}{\textwidth}{@{}P{1.6cm}YP{0.6cm}P{0.6cm}P{1.0cm}@{}}
\toprule
\textbf{Tarefa / versão} & \textbf{Sequência (M = preparação; T = toque = P+K)} & \textbf{$n_M$} & \textbf{$A_{\min}$} & \textbf{Tempo (s)} \\
\midrule
C1-A & M T(interruptor) M T(Histórico) M T(Concluir) M T(causa) T(Confirmar) & 4 & 5 & 12,30 \\
C1-B & M T(``Por quê?'') M T(Desfazer) M T(Concluir) M T(causa) T(Confirmar) & 4 & 5 & 12,30 \\
C2 (pausar) & M T(Rotinas) M T(``Bom dia'') M T(pausa) M T(Concluir) M T(Confirmar) & 5 & 5 & 13,65 \\
C2 (editar) & M T(Rotinas) M T(``Bom dia'') M T(editar) M T(Cafeteira) T(Salvar) M T(Concluir) M T(Confirmar) & 6 & 7 & 17,76 \\
V1-A & M T(Histórico) M T(Rotinas) M T(``Horário de pico'') M T(pausa) M T(Concluir) M T(causa) T(Confirmar) & 6 & 7 & 17,76 \\
V1-B & M T(``Por quê?'') M T(``Ir para a rotina'') M T(pausa) M T(Concluir) M T(causa) T(Confirmar) & 5 & 6 & 15,03 \\
V2-A & M T(Histórico) M T(Rotinas) M T(``Climatização'') M T(editar) M 6T($+$) M T(Salvar) M T(Concluir) M T(causa) T(Confirmar) & 8 & 14 & 30,12 \\
V2-B & M T(``Por quê?'') M T(``Ir para a rotina'') M T(editar) M 6T($+$) M T(Salvar) M T(Concluir) M T(causa) T(Confirmar) & 7 & 13 & 27,39 \\
\bottomrule
\end{tabularx}

\end{table}

Tempo = $n_M \times 1{,}35 + A_{\min} \times 1{,}38$. $A_{\min}$ é o denominador de $X = A/A_{\min}$. Os valores de V1 e V2 servem só de denominador e de referência; não foram usados para escolher nenhuma alternativa de design.

\subsection{Comparação de duas alternativas para a cena principal}
A cena principal é o diagnóstico da causa na versão~B de C1, onde se concentra a pergunta do estudo. Alternativas: \textbf{B1} explicação sob demanda (o aviso tem ``Por quê?''; a causa só aparece no painel após o toque) e \textbf{B2} explicação sempre exibida (o aviso já traz a causa, como na condição de Li \textit{et al.}~\cite{li2026}).

\begin{table}[!h]
\caption{KLM das alternativas da cena principal (C1, versão B)}
\label{tab:klm-alt}
\centering
\footnotesize
\begin{tabularx}{\textwidth}{@{}P{3.2cm}YP{0.6cm}P{0.6cm}P{1.0cm}@{}}
\toprule
\textbf{Alternativa} & \textbf{Sequência} & \textbf{$n_M$} & \textbf{T} & \textbf{Tempo (s)} \\
\midrule
B1 Sob demanda (adotada) & M T(``Por quê?'') M T(Desfazer) M T(Concluir) M T(causa) T(Confirmar) & 4 & 5 & 12,30 \\
B2 Sempre exibida & M T(Desfazer) M T(Concluir) M T(causa) T(Confirmar) & 3 & 4 & 9,57 \\
\bottomrule
\end{tabularx}
\end{table}

\textit{Decisão (D4).} B2 é 2,73~s e um toque mais rápida para um especialista, mas B1 é adotada porque: (a)~a pergunta de pesquisa é sobre explicação \textit{sob demanda}; (b)~em B1 a leitura da explicação é observável (\ev{support\_open(por\_que)}), enquanto em B2 ela seria uma operação não observável, e o uso de suporte $U$ ficaria indefinido na versão~B; (c)~B2 acrescenta duas linhas de texto a uma tela já densa, com mais alvos próximos aos interruptores (risco de E-SLP, \stHO). O KLM descreve um especialista sem erros e não decide sozinho; ele quantifica o custo aceito. Perguntas que a comparação levantou: o painel deve fechar sozinho após ``Desfazer''? (Não: a pessoa confere o resultado; fecha por ``Fechar'' ou pela navegação.) ``Desfazer'' deve ficar também no aviso? (Sim, para que B não dependa de abrir o painel para restaurar; isso mantém $A_{\min}$ igual a A.)

\subsection{Escopo do GOMS neste trabalho}
O GOMS descreve usuários experientes, sem erros, em tarefas rotineiras~\cite{barbosa2010}, \cite{preece2002}. Aqui ele serve para: (1)~comparar alternativas de design; (2)~definir o caminho de referência e o denominador $A_{\min}$; (3)~detectar métodos inconsistentes para objetivos semelhantes (pausar em C2 e V1 usa o mesmo método, o que confirma a consistência de S3); e (4)~gerar o conteúdo da ajuda procedimental (Seção~\ref{sec:ajuda}). Ele \textit{não} é modelo do comportamento da persona: usar os tempos do KLM como tempo de uma persona novata seria hipótese sem sustentação, e não há tempo comportamental simulado no UXAgent. \textit{Fronteira:} HTA, CTT, métodos, regras de seleção e $A_{\min}$ são modelos do analista e não são entregues ao agente.

% =====================================================================
\section{Interação como conversa: MoLIC}\label{sec:molic}

\subsection{Diagrama}
As Figuras~\ref{fig:molic-a} e~\ref{fig:molic-b} trazem a MoLIC do único papel de usuário (a moradora)~\cite{barbosa2003molic}, \cite{barbosa2010}. Vocabulário controlado: verbo no infinitivo + objeto (``abrir rotina'', ``pausar rotina''); falas do usuário (u:) são exatamente as ações que o agente pode escolher naquela cena; falas do preposto (d:) são o que o agente observará em seguida. Falas de recuperação de ruptura estão tracejadas.

\begin{sidewaysfigure}
\centering
\resizebox{\textheight}{!}{%
\begin{tikzpicture}[
  scene/.style={draw, rounded corners=6pt, text width=3.3cm, align=left, font=\scriptsize, inner sep=3pt, fill=white},
  proc/.style={fill=black, minimum width=6mm, minimum height=3mm, inner sep=0pt},
  ua/.style={draw, fill=gray!25, text width=7.6cm, align=center, font=\scriptsize, inner sep=4pt},
  tr/.style={-{Latex[length=1.8mm]}, thin},
  rec/.style={tr, dashed},
  lab/.style={font=\tiny, align=center, fill=white, inner sep=1pt},
  conn/.style={font=\scriptsize\itshape, align=center}
]
\node[ua] (ua) at (9.2,3.6) {\textbf{Acesso ubíquo}\\barra inferior: Início, Rotinas, Histórico, Ajuda\\barra superior: Concluir, Desistir};
\node[conn] (toS2) at (9.2,5.6) {S2 Rotinas (Fig.~\ref{fig:molic-b})};
\node[circle, fill=black, minimum size=3mm, inner sep=0pt] (open) at (-3.0,0) {};
\node[scene] (s1) at (0,0) {\cena{S1 Início}{Ver a casa}{d: aparelhos e estados\\d: aviso de mudança\\u: alternar aparelho\\u: ver por quê \textit{[B]}\\u: desfazer mudança \textit{[B]}}};
\node[scene] (s8) at (4.8,0) {\cena{S8 Painel ``Por quê?'' [B]}{Entender a ação automática}{d: causa (rotina, hora, condição)\\u: desfazer mudança\\u: ir para a rotina\\u: fechar painel}};
\node[scene] (s5) at (9.6,0) {\cena{S5 Histórico}{Saber quem mudou}{d: hora, aparelho, novo estado, origem}};
\node[scene] (s6) at (14.4,0) {\cena{S6 Ajuda}{Tirar uma dúvida}{d: perguntas (Tab.~\ref{tab:ajuda})\\u: abrir tópico\\u: voltar}};
\node[scene] (s7) at (19.2,0) {\cena{S7 Concluir}{Encerrar a tarefa}{d: Por que isso aconteceu?\\u: escolher resposta\\u: confirmar\\u: voltar}};
\node[proc, label={[font=\tiny]left:alternar aparelho}] (p1) at (-1.0,-3.0) {};
\node[proc, label={[font=\tiny]below:desfazer mudança}] (p2) at (2.4,-3.0) {};
\node[conn] (toS3) at (6.6,-3.0) {S3 Rotina (Fig.~\ref{fig:molic-b})};
\node[conn] (orig6) at (14.4,-3.0) {cena de origem};
\node[proc, label={[font=\tiny]left:registrar fim}] (p8) at (18.4,-3.0) {};
\node[proc, label={[font=\tiny]right:registrar desistência}] (p9) at (27.4,-3.0) {};
\node[conn] (orig7) at (21.6,-3.0) {cena de origem};
\node[scene, double, text width=2.6cm] (ce2) at (24.6,0) {\cena{Captura (CE2)}{Resposta ausente}{d: Escolha uma resposta antes de confirmar.}};
\node[circle, draw, double, double distance=1pt, fill=black, minimum size=3mm, inner sep=0pt] (close) at (22.8,-5.2) {};
\node[font=\tiny, align=left, anchor=west] at (23.1,-5.2) {fechamento:\\verificador externo\\lê o estado};
% abertura
\draw[tr] (open) -- node[lab, above]{d: instrução\\da tarefa} (s1);
% S1
\draw[tr] ([xshift=-6mm]s1.south) -- node[lab, left]{u: alternar\\aparelho} (p1);
\draw[tr] (p1) -- node[lab, right, pos=0.6]{d: novo\\estado} ([xshift=-2mm]s1.south);
\draw[tr] ([xshift=6mm]s1.south) -- node[lab, left]{u: desfazer [B]} (p2);
\draw[tr] ([yshift=2mm]s1.east) -- node[lab, above]{u: ver por quê [B]} ([yshift=2mm]s8.west);
\draw[tr] ([yshift=-4mm]s8.west) -- node[lab, below]{u: fechar} ([yshift=-4mm]s1.east);
\draw[tr] ([xshift=-6mm]s8.south) -- node[lab, right]{u: desfazer} (p2);
\draw[tr] (p2) to[bend left=25] node[lab, below]{d: aparelho restaurado} ([xshift=2mm]s1.south);
\draw[tr] ([xshift=6mm]s8.south) -- node[lab, right]{u: ir para a rotina} (toS3);
% S6 retorno
\draw[tr] (s6) -- node[lab, right]{u: voltar\\d: retorno} (orig6);
% S7
\draw[tr] ([xshift=-4mm]s7.south) -- node[lab, left]{u: confirmar\\(com resposta)} (p8);
\draw[tr] (p8) -- node[lab, left]{d: tarefa encerrada} (close);
\draw[tr] ([yshift=3mm]s7.east) -- node[lab, above]{u: confirmar\\sem resposta} ([yshift=3mm]ce2.west);
\draw[rec] ([yshift=-4mm]ce2.west) -- node[lab, below]{u: escolher\\resposta} ([yshift=-4mm]s7.east);
\draw[tr] ([xshift=12mm]s7.south) -- node[lab, right]{u: voltar} (orig7);
\draw[tr] (p9) -- node[lab, right]{d: tarefa\\interrompida} (close);
% acesso ubíquo
\draw[tr] (ua.south west) -- node[lab, above left]{u: ir para Início} (s1.north);
\draw[tr] ([xshift=4mm]ua.south) -- node[lab, right]{u: ver histórico} (s5.north);
\draw[tr] ([xshift=28mm]ua.south) -- node[lab, right, pos=0.4]{u: pedir ajuda} (s6.north);
\draw[tr] (ua.south east) -- node[lab, above right]{u: concluir} (s7.north);
\draw[tr] (ua.north) -- node[lab, right]{u: ver rotinas} (toS2);
\draw[tr] (ua.east) -| node[lab, above, pos=0.25]{u: desistir} (p9);
\end{tikzpicture}}

\caption{MoLIC, parte (a): navegação, diagnóstico e encerramento. Retângulos arredondados: cenas (tópico e diálogos); retângulos pretos: processos do sistema; cinza: acesso ubíquo, disponível em todas as cenas; tracejado: fala de recuperação de ruptura.}
\label{fig:molic-a}
\end{sidewaysfigure}

\begin{figure}[!t]
\centering
\resizebox{\textwidth}{!}{%
\begin{tikzpicture}[
  scene/.style={draw, rounded corners=6pt, text width=3.3cm, align=left, font=\scriptsize, inner sep=3pt, fill=white},
  alert/.style={scene, double},
  proc/.style={fill=black, minimum width=6mm, minimum height=3mm, inner sep=0pt},
  tr/.style={-{Latex[length=1.8mm]}, thin},
  rec/.style={tr, dashed},
  lab/.style={font=\tiny, align=center, fill=white, inner sep=1pt},
  conn/.style={font=\scriptsize\itshape, align=center}
]
\node[conn] (inS2) at (0,3.0) {acesso ubíquo:\\ver rotinas};
\node[conn] (inS3) at (3.4,3.0) {S8: ir para\\a rotina [B]};
\node[scene] (s2) at (0,0) {\cena{S2 Rotinas}{Ver as rotinas}{d: nome e situação\\u: abrir rotina\\u: desfazer exclusão (faixa)}};
\node[scene] (s3) at (5.2,0) {\cena{S3 Rotina}{Ver e controlar a rotina}{d: regra se\dots então, situação\\u: pausar/retomar (ícone)\\u: editar rotina (ícone)\\u: excluir rotina\\u: ver dica (?)\\u: voltar}};
\node[alert] (al) at (10.8,0) {\cena{Alerta (AL)}{Confirmar exclusão}{d: Excluir a rotina X? Ela sai da lista.\\u: excluir\\u: cancelar}};
\node[proc, label={[font=\tiny]left:restaurar rotina}] (p5) at (-1.0,-3.2) {};
\node[proc, label={[font=\tiny]above:pausar/retomar}] (p3) at (7.4,3.0) {};
\node[proc, label={[font=\tiny]right:excluir rotina}] (p4) at (10.8,3.0) {};
\node[scene] (s4) at (5.2,-6.0) {\cena{S4 Editar rotina}{Mudar horário e aparelhos}{d: Smart Schedule hh:00 (00h--23h)\\d: aparelhos com caixas\\u: aumentar/diminuir horário\\u: marcar/desmarcar aparelho\\u: salvar \quad u: cancelar\\u: ver dica (?)}};
\node[proc, label={[font=\tiny]right:salvar rotina}] (p7) at (8.0,-3.0) {};
\node[proc, label={[font=\tiny]right:ajustar valor (PA: 00h--23h)}] (p6) at (5.2,-9.6) {};
\node[alert] (ce) at (11.2,-6.0) {\cena{Captura (CE1)}{Alterações não salvas}{d: Você mudou esta rotina. Salvar as mudanças?\\u: salvar\\u: sair sem salvar}};
\node[scene] (s9) at (0,-6.0) {\cena{S9 Dica local}{Explicar os símbolos da tela}{d: significado dos ícones e de Smart Schedule\\u: fechar}};
\node[conn] (orig9) at (0,-9.6) {cena de origem};
% entradas
\draw[tr] (inS2) -- (s2);
\draw[tr] (inS3) -- ([xshift=-12mm]s3.north);
% S2
\draw[tr] ([yshift=3mm]s2.east) -- node[lab, above]{u: abrir rotina} ([yshift=3mm]s3.west);
\draw[tr] ([yshift=-8mm]s3.west) -- node[lab, below]{u: voltar} ([yshift=-8mm]s2.east);
\draw[rec] ([xshift=-8mm]s2.south) -- node[lab, left]{u: desfazer\\exclusão (RA)} (p5);
\draw[tr] (p5) -- node[lab, right]{d: rotina\\de volta} ([xshift=-1mm]s2.south);
% S3
\draw[tr] ([xshift=10mm]s3.north) -- node[lab, left]{u: pausar|retomar} (p3);
\draw[tr] (p3) -- node[lab, right]{d: situação\\atualizada} ([xshift=17mm]s3.north east);
\draw[tr] ([yshift=4mm]s3.east) -- node[lab, above]{u: excluir rotina} ([yshift=4mm]al.west);
\draw[rec] ([yshift=-6mm]al.west) -- node[lab, below]{u: cancelar} ([yshift=-6mm]s3.east);
\draw[tr] (al) -- node[lab, right]{u: excluir} (p4);
\draw[tr] (p4.north) -- ++(0,0.8) -| node[lab, above, pos=0.25]{d: rotina excluída + faixa ``Desfazer''} ([xshift=-14mm]s2.north);
\draw[tr] ([xshift=-6mm]s3.south) -- node[lab, left]{u: editar rotina} ([xshift=-6mm]s4.north);
\draw[tr] ([xshift=-14mm]s3.south) -- node[lab, above left]{u: ver dica} (s9.north east);
% S4
\draw[tr] ([xshift=10mm]s4.north) -- node[lab, left, pos=0.3]{u: salvar} (p7);
\draw[tr] (p7) -- node[lab, right]{d: rotina salva} ([xshift=12mm]s3.south);
\draw[tr] ([xshift=-4mm]s4.south) -- node[lab, left]{u: aumentar|diminuir|\\marcar|desmarcar} (p6);
\draw[tr] (p6) -- node[lab, right]{d: valor exibido} ([xshift=4mm]s4.south);
\draw[tr] (s4.east) -- node[lab, above]{u: cancelar} node[lab, below]{(com alterações)} (ce.west);
\draw[tr] (ce.north west) -- node[lab, right]{u: salvar} (p7);
\draw[rec] (ce.north) to[out=95, in=-20] node[lab, right, pos=0.5]{u: sair sem salvar} ([xshift=16mm]s3.south east);
\draw[tr] (s4.west) -- node[lab, above]{u: ver dica} (s9.east);
\draw[tr] (s9) -- node[lab, right]{u: fechar\\d: retorno} (orig9);
\end{tikzpicture}}

\caption{MoLIC, parte (b): gestão de rotinas. Borda dupla: cenas de alerta (AL) e de captura de erro (CE). ``Cancelar'' sem alterações volta direto a S3 (omitido). O acesso ubíquo da Fig.~\ref{fig:molic-a} vale em todas as cenas.}
\label{fig:molic-b}
\end{figure}

\subsection{Leituras de cada elemento}
\begin{table}[!h]
\caption{Leituras de cada elemento da MoLIC}
\label{tab:leituras}
\centering
\footnotesize
\begin{tabularx}{\textwidth}{@{}P{2.4cm}YYP{3.6cm}@{}}
\toprule
\textbf{Elemento} & \textbf{Leitura da persona} & \textbf{Leitura do agente (observação ou ação)} & \textbf{Evento registrado} \\
\midrule
Cena & Uma tela com um assunto (``ver a casa'', ``ver a rotina'') & Observação $o_t$ com título, signos e lista de ações & \ev{screen\_view} (tela, elementos visíveis) \\
Fala u: & Algo que ela pode tocar ou alternar ali & Item numerado da lista de ações; é a única coisa que o agente pode responder & \ev{action\_proposed}$^*$, \ev{action\_executed} \\
Fala d: & O que a tela mostra depois do toque & Próxima $o'_t$, com a linha ``Mudou na tela'' quando for o caso & \ev{screen\_view} ou \ev{state\_changed} \\
Processo do sistema & O aparelho ou a rotina muda & Invisível como processo; só o resultado aparece na próxima observação & \ev{state\_changed} (variáveis, origem) \\
Fala de recuperação & Desfazer, cancelar, voltar & Ação comum da lista; não é marcada como ``recuperação'' para o agente & \ev{action\_executed}; \ev{dialog\_choice}; etiqueta MAC derivada \\
Alerta e captura & Uma caixa que pergunta ou avisa & Observação da caixa com suas duas opções & \ev{dialog\_shown} (AL ou CE), \ev{dialog\_choice} \\
Acesso ubíquo à ajuda & Aba Ajuda e botão ``?'' em qualquer tela & Itens da lista em toda observação (``tocar Ajuda'', ``tocar ?''); o texto só chega após o toque & \ev{support\_open}, \ev{support\_close} \\
\bottomrule
\end{tabularx}

\end{table}
\noindent{\scriptsize $^*$ Só sintético.}

\subsection{Da MoLIC à máquina de estados}
A MoLIC é ferramenta epistêmica e não modelo formal processável~\cite{barbosa2010}; a Tabela~\ref{tab:estados} a traduz numa rede de transição de estados que o executor processa~\cite[cap.~16]{dix2004}. Regras: as falas u: de um estado são as ações disponíveis ao executor naquele estado e o conteúdo da observação seguinte; a conclusão de um objetivo é uma condição do estado da interface avaliada pelo verificador externo, nunca pela declaração do agente; a ajuda é subdiálogo que guarda a cena de origem numa pilha e devolve a persona a ela.

\begin{table}[!h]
\caption{Máquina de estados da interação (UA: acesso ubíquo, disponível em todos os estados exceto diálogos modais)}
\label{tab:estados}
\centering
\scriptsize
\begin{tabularx}{\textwidth}{@{}P{1.4cm}P{3.8cm}YP{1.8cm}P{3.2cm}@{}}
\toprule
\textbf{Estado} & \textbf{Falas u: (ações oferecidas)} & \textbf{Processo / fala d:} & \textbf{Próximo} & \textbf{Eventos} \\
\midrule
UA & Início; Rotinas; Histórico; Ajuda; Concluir; Desistir & Ajuda e Concluir empilham a cena de origem; Desistir encerra & S1; S2; S5; S6; S7; FIM & \ev{screen\_view}; \ev{support\_open}; \ev{task\_end(desistencia)} \\
S1 & alternar <aparelho> (4); [B] Por quê?; [B] Desfazer & alternar: novo estado; Desfazer: restaura o aparelho, aviso some & S1; S8; S1 & \ev{action\_executed}, \ev{state\_changed(origem=pessoa)}; \ev{support\_open(por\_que)} \\
S8 [B] & Desfazer; Ir para a rotina; Fechar & Desfazer: restaura; painel mostra ``Feito'' & S8; S3; S1 & \ev{state\_changed}; \ev{screen\_view}; \ev{support\_close} \\
S2 & abrir <rotina>; Desfazer exclusão (se faixa visível) & restaurar rotina excluída & S3; S2 & \ev{screen\_view}; \ev{state\_changed(restaurada)} \\
S3 & Voltar; pausar|retomar (ícone); editar (ícone); ?; Excluir rotina & pausar/retomar: situação alterna & S2; S3; S4; S9; AL & \ev{state\_changed(status)}; \ev{support\_open(dica)}; \ev{dialog\_shown(AL)} \\
AL & Cancelar; Excluir & excluir: rotina sai da lista; faixa ``Desfazer'' em S2 & S3; S2 & \ev{dialog\_choice}; \ev{state\_changed(excluida)} \\
S4 & Cancelar; $-$; $+$; marcar|desmarcar <aparelho>; ?; Salvar & $\pm$: horário $\pm$1~h, preso em 00h--23h (PA); Salvar: grava; Cancelar com alterações: CE1 & S4; S3; CE1; S9 & \ev{action\_executed}; \ev{state\_changed}; \ev{dialog\_shown(CE)} \\
CE1 & Salvar; Sair sem salvar & salvar grava; sair descarta & S3 & \ev{dialog\_choice}; \ev{state\_changed} \\
S5 & (só UA) & --- & --- & \ev{screen\_view(S5)} \\
S6 & abrir <tópico>; Voltar & mostra a resposta; Voltar desempilha & S6; origem & \ev{support\_open(ajuda, tópico)}; \ev{support\_close} \\
S9 & Fechar & desempilha & origem & \ev{support\_close} \\
S7 & escolher <opção>; Confirmar; Voltar & Confirmar com opção: registrar fim; sem opção: CE2 & S7; FIM; S7; origem & \ev{cause\_selected}; \ev{task\_end(concluida)}; \ev{dialog\_shown(CE)} \\
FIM & --- & verificador externo lê o JSON final e decide sucesso & --- & \ev{task\_end}; \ev{error\_tagged}; \ev{method\_selected}; \ev{communicability\_tagged} \\
(controlador) & --- & limite de $3A_{\min}$ ações ou falha técnica & FIM & \ev{task\_end(limite|indeterminado)}; \ev{technical\_failure} \\
\bottomrule
\end{tabularx}

\end{table}

% =====================================================================
\section{Rupturas: prevenção e recuperação}\label{sec:rupturas}

Para cada erro previsto nas HTAs da Fase~II escolheu-se o mecanismo~\cite{barbosa2010}: \textbf{PP} prevenção passiva (o signo ou a documentação informa); \textbf{PA} prevenção ativa (o sistema impede o erro); \textbf{AL} prevenção apoiada (o sistema alerta e pede decisão); \textbf{RA} recuperação apoiada (o sistema oferece desfazer ou corrigir); \textbf{CE} captura de erro (o sistema detecta e avisa). Regras verificadas: toda PP tem RA equivalente, e toda PA sobre signo de operação vem com PP.

\begin{table}[!h]
\caption{Matriz de rupturas por operação da HTA}
\label{tab:rupturas}
\centering
\scriptsize
\setlength{\tabcolsep}{3pt}
\begin{tabularx}{\textwidth}{@{}P{2.3cm}P{2.3cm}YP{2.0cm}P{2.3cm}P{1.3cm}@{}}
\toprule
\textbf{Operação da HTA} & \textbf{Erro previsto} & \textbf{Mecanismo} & \textbf{Signo} & \textbf{Evento} & \textbf{Métrica} \\
\midrule
C1.1.2 Tocar Histórico; V1.1, V2.1 & Engano: não busca a origem (``Cadê?'') & PP: aba ``Histórico'' com texto em toda tela; tópico H5 da Ajuda. RA: a busca pode ser feita a qualquer momento antes de Confirmar & barra.navegação & \ev{screen\_view(S5)} & $R_{\mathrm{ATR}}$, $U$ \\
C1.1.4 e 3.2 Decidir e escolher a causa & E-ATR & PP: origem escrita em S5; [B] painel ``Por quê?''. RA: ``Voltar'' em S7 e troca de opção antes de Confirmar & historico.linha; painel.causa & \ev{cause\_selected} & $R_{\mathrm{ATR}}$ \\
C1.2.3 Tocar interruptor & E-SLP (aciona o vizinho) & PP: nome e estado em texto ao lado de cada interruptor. RA: tocar de novo; aviso e Histórico mostram a mudança & aparelho.estado & \ev{action\_executed}, \ev{state\_changed} & $R_{\mathrm{SLP}}$, $X$ \\
C1.3.1, C2.3 Tocar Concluir & E-OMI (conclui sem o estado-alvo) & PP: instrução do protocolo. RA: ``Voltar'' em S7. Sem PA nem CE: o sistema não conhece o objetivo da tarefa, e avisar seria entregar o critério de sucesso & barra.concluir & \ev{task\_end} & $S$, $R_{\mathrm{OMI}}$ \\
C1.3.3 Confirmar & Confirmar sem resposta & CE2: ``Escolha uma resposta antes de confirmar.'' PP: pergunta sem opção pré-marcada & concluir.pergunta & \ev{dialog\_shown(CE)} & --- (falha de operação, não erro de tarefa) \\
C2.1.2 / V1.2.3 Escolher a rotina & E-ACT (age na rotina errada) & PP: nome e situação em S2; regra completa em S3 antes de agir. RA: retomar (mesmo botão) & rotina.nome; rotina.regra & \ev{state\_changed(status)} & $R_{\mathrm{ACT}}$ \\
C2.2.1 / V1.3.1 Tocar ícone de pausa & Engano: não reconhece o ícone (``O que é isto?'') & PP: dica local ``?'' e tópico H2. RA: o mesmo botão retoma & rotina.pausar & \ev{support\_open}, \ev{state\_changed} & $U$, $X$ \\
C2.2.1 / V1.3.1 (adjacente) & E-ACT: toca ``Excluir rotina'' & PP: texto vermelho ``Excluir rotina'', separado dos ícones. AL: alerta de confirmação. RA: faixa ``Desfazer'' em S2 & rotina.excluir; alerta & \ev{dialog\_shown(AL)}, \ev{dialog\_choice}, \ev{state\_changed} & $R_{\mathrm{ACT}}$ \\
C2.2.2 Tocar ícone de editar & D: toca a pausa (adjacente) & PP: dica local. RA: retomar & rotina.editar & \ev{action\_executed} & $R_{\mathrm{SLP}}$ \\
C2.2.3 Desmarcar aparelho & E-SLP (caixa vizinha) & PP: nome ao lado da caixa. RA: marcar de novo ou ``Cancelar'' & rotina.aparelhos & \ev{action\_executed} & $R_{\mathrm{SLP}}$ \\
C2.2.4 / V2.4 Tocar Salvar & E-OMI: sai sem salvar & PP: botão ``Salvar'' em texto. CE1: ``Você mudou esta rotina. Salvar as mudanças?'' RA: ``Salvar'' no próprio CE1 & captura.CE1 & \ev{dialog\_shown(CE)}, \ev{dialog\_choice} & $R_{\mathrm{OMI}}$ \\
V2.3.1 Tocar ``$+$'' & E-SLP: toca ``$-$'' & PP: sinais grandes e valor em destaque. RA: o botão inverso & rotina.horário & \ev{action\_executed} & $R_{\mathrm{SLP}}$, $X$ \\
V2.3.1 (limite) & Passar de 23h para 0h & PA: valor preso em 00h--23h. PP: nota ``de 00h a 23h''. RA: ``$-$'' & rotina.horário & \ev{action\_executed} sem mudança & --- \\
V2.3.2 Conferir horário & L: perde a conta & PP: horário atualizado após cada toque (signo dinâmico). RA: ``$+$''/``$-$'' & rotina.horário & --- (inferido) & $X$ \\
\bottomrule
\end{tabularx}

\end{table}

\textit{O mecanismo muda a métrica.} Com PA o erro não pode ocorrer (passar de 23h para 0h não é erro de tarefa). Com PP e RA ele ocorre e é registrado. Com AL a ação destrutiva só conta como E-ACT quando a exclusão é confirmada; abrir o alerta e cancelar é registrado como \ev{dialog\_choice(cancelar)} e etiquetado ``Epa!'' (Seção~\ref{sec:mac}). \textit{Restrição do comparador.} Pessoas idosas e o grupo de 30 a 50 anos usarão o mesmo CasaFácil no piloto, com os mesmos mecanismos; por isso as taxas de erro sintéticas e humanas medem a mesma coisa. A comparação com aplicativos comerciais, cuja interface não é conhecida em detalhe, é limitação declarada.

% =====================================================================
\section{Comunicabilidade operacionalizada (MAC)}\label{sec:mac}

Oito etiquetas do MAC~\cite{prates2000}, \cite{barbosa2010} são definidas como padrões de eventos no registro (Tabela~\ref{tab:mac}). Não se infere ruptura do texto da justificativa do LLM. O verificador grava cada ocorrência como \ev{communicability\_tagged} com a etiqueta e os eventos de origem. Cada definição é \stHO: o padrão no registro do agente não equivale à experiência de ruptura de uma pessoa. A mesma definição será aplicada aos registros humanos do piloto.

\begin{table}[!h]
\caption{Etiquetas do MAC como padrões de eventos}
\label{tab:mac}
\centering
\scriptsize
\begin{tabularx}{\textwidth}{@{}P{2.2cm}YP{4.4cm}@{}}
\toprule
\textbf{Etiqueta} & \textbf{Assinatura no registro} & \textbf{Onde se espera e o que indica} \\
\midrule
``Cadê?'' & Três ou mais \ev{screen\_view} seguidos por acesso ubíquo ou Voltar, sem \ev{state\_changed} nem \ev{support\_open} entre eles & C1/V1/V2 em A, procurando a origem; C2 procurando as rotinas. Signo de navegação que não comunica onde está a função \\
``O que é isto?'' & \ev{support\_open} na camada dica ou nos tópicos H2, H3 ou H4 & S3 e S4: ícones sem texto e ``Smart Schedule'' (L2) \\
``Epa!'' & \ev{action\_executed} seguido, em até duas ações, da ação inversa sobre o mesmo elemento (alternar de volta, $-$ após $+$, retomar após pausar, Desfazer exclusão) ou \ev{dialog\_choice(cancelar)} no AL & Interruptores adjacentes, ícones adjacentes, $\pm$ (L3, L4) \\
``Onde estou?'' & \ev{screen\_view(X)} seguido de saída de X na ação seguinte, sem nenhuma outra ação em X, duas ou mais vezes na tarefa & Entra e sai de telas sem agir; cena não se identifica \\
``Por que não funciona?'' & A mesma ação executada duas ou mais vezes seguidas sem \ev{state\_changed} (por exemplo, Confirmar sem resposta; $+$ em 23h) & S7 e S4; o sistema não responde como esperado \\
``Socorro!'' & \ev{support\_open} na camada ajuda (S6), qualquer tópico & Qualquer cena; complementa ``O que é isto?'' \\
``Para mim está bom'' & \ev{task\_end(concluida)} e verificador = fracasso & E-OMI e E-ATR: a pessoa acha que terminou \\
``Desisto'' & \ev{task\_end(desistencia)} & Qualquer cena; ruptura total \\
\bottomrule
\end{tabularx}

\end{table}

% =====================================================================
\section{Sistema de ajuda}\label{sec:ajuda}

A ajuda é a metacomunicação explícita do designer~\cite{barbosa2010}, \cite{silveira2002} e, neste estudo, também variável de protocolo: a ajuda permitida (Fase~II) e o uso de suporte, uma das métricas mínimas. Ela é igual para pessoas e para o agente, não é adaptativa e só chega ao agente depois de uma ação registrada de abri-la (\textit{PedirAjuda})~\cite{lara2026p3}.

\subsection{Conteúdo a partir das dúvidas}
\begin{table}[!h]
\caption{Tópicos da Ajuda (S6) por tipo de dúvida; texto igual ao exibido}
\label{tab:ajuda}
\centering
\scriptsize
\begin{tabularx}{\textwidth}{@{}P{0.5cm}P{1.6cm}P{3.0cm}YP{1.8cm}@{}}
\toprule
\textbf{Id} & \textbf{Tipo de dúvida} & \textbf{Pergunta (como aparece)} & \textbf{Resposta} & \textbf{Modelo de origem} \\
\midrule
H1 & Informativa; motivacional & O que é uma rotina? & ``Uma rotina é uma regra que o aplicativo segue sozinho: num horário ou numa situação, ela liga ou desliga aparelhos. Quem cria a rotina escolhe quando ela age. As rotinas ficam na aba Rotinas.'' & Esquema de signos \\
H2 & Descritiva; interpretativa & O que faz o botão de duas barras? & ``É o botão Pausar. Uma rotina pausada fica guardada, mas não age até ser retomada. Quando ela está pausada, o botão mostra um triângulo: é o Retomar.'' & Esquema de signos \\
H3 & Descritiva & O que faz o botão de lápis? & ``É o botão Editar. Ele abre a rotina para mudar o horário e os aparelhos incluídos.'' & Esquema de signos \\
H4 & Interpretativa & O que é ``Smart Schedule''? & ``É o horário em que a rotina age. Toque em $+$ ou $-$ para mudar de hora em hora, de 00h a 23h. Depois, toque em Salvar.'' & Esquema de signos; CTT \\
H5 & Investigativa; histórica & Como saber quem mudou um aparelho? & ``Abra a aba Histórico. Cada linha mostra a hora, o aparelho, o que mudou e quem mudou: `Você' ou o nome da rotina.'' \textit{[B: ``Você também pode tocar em `Por quê?' no aviso da tela Início.'']} & MoLIC (S5, S8) \\
H6 & Procedimental & Como fazer uma rotina parar de agir por um tempo? & ``Abra Rotinas, toque no nome da rotina e toque no botão de duas barras (Pausar). Para ela voltar a agir, toque no triângulo (Retomar).'' & CMN-GOMS (método PAUSAR) \\
H7 & De escolha & Pausar, editar ou excluir: qual usar? & ``Pausar: a rotina para de agir, mas fica guardada. Editar: muda o horário ou os aparelhos. Excluir: apaga a rotina; o aplicativo pede confirmação e, logo depois, ainda mostra `Desfazer' na lista de rotinas.'' & GOMS (regras de seleção); matriz de rupturas \\
H8 & Navegacional & Onde ficam as rotinas? & ``Na aba Rotinas, na barra de baixo, que aparece em todas as telas.'' & MoLIC (acesso ubíquo) \\
H9 & Procedimental (recuperação) & Como desfazer uma mudança? & ``Se você ligou ou desligou um aparelho sem querer, toque no interruptor de novo. Se pausou uma rotina, toque em Retomar. Se excluiu uma rotina, toque em Desfazer na lista de rotinas.'' \textit{[B: ``Se uma rotina mudou um aparelho, toque em Desfazer no aviso.'']} & Matriz de rupturas (RA) \\
H10 & Informativa (protocolo) & O que são Concluir e Desistir? & ``Concluir: avisa que você terminou o que foi pedido. Desistir: encerra sem terminar.'' & Protocolo \\
\bottomrule
\end{tabularx}

\end{table}

As respostas usam os mesmos termos da tela e são genéricas: nenhuma menciona aparelhos, rotinas ou horários de uma tarefa específica. Os tópicos foram derivados do esquema de signos, da MoLIC, do CMN-GOMS de C2 e da matriz de rupturas, e não das tarefas V1 e V2. A diferença entre A e B se limita às frases que mencionam o painel e o ``Desfazer'' do aviso, que só existem em B; ela é diferença entre versões, e não entre pessoas e agente. Dúvidas sugestivas (``o que mais posso fazer?'') não têm tópico próprio, porque o aplicativo tem poucas funções e H1, H7 e H8 as cobrem; essa ausência é decisão, não descuido.

\subsection{Acesso e apresentação}
Três camadas, todas sob demanda e contextuais~\cite{barbosa2010}:
\begin{enumerate}
\item \textbf{Dica local} (S9): botão ``?'' em S3 e S4, as duas telas com signos de aplicação sem texto. Abre um balão sobre a própria tela com H2, H3 e H4 resumidos, sem sair da cena.
\item \textbf{Ajuda} (S6): aba da barra inferior, acessível em todas as cenas (acesso ubíquo); lista de perguntas da Tabela~\ref{tab:ajuda}; ``Voltar'' devolve à cena de origem.
\item \textbf{Explicação da ação automática} (S8, só versão B): painel ``Por quê?'' do aviso. É a condição experimental, mas é registrada como suporte (\texttt{por\_que}) para que $U$ seja comparável entre versões.
\end{enumerate}
O caderno da persona não é camada do sistema: é artefato próprio, registrado só no agente (M5, \texttt{caderno}).

\textit{Requisitos de Dix}~\cite[cap.~11]{dix2004}. \textit{Disponibilidade}: Ajuda em todas as cenas; dica local onde há signos sem texto. \textit{Precisão e completude}: cada signo de aplicação das Tabelas~\ref{tab:signos1} e~\ref{tab:signos2} tem tópico; conferido por lista de verificação. \textit{Consistência}: mesmo vocabulário dos rótulos; a ajuda chama o ícone pela forma, como a persona o vê. \textit{Robustez}: abrir a ajuda nunca muda o estado dos aparelhos nem das rotinas, e o retorno restaura a cena. \textit{Flexibilidade}: duas profundidades (dica e Ajuda) e perguntas na linguagem da persona; não há busca por texto, porque o protocolo não tem digitação (ausência justificada). \textit{Discrição}: a ajuda nunca abre sozinha e não interrompe a tarefa; as lições do Clippy sobre intrusão e controle valem aqui~\cite{dix2004}, \cite{preece2002}.

\subsection{Registro e métrica}
Cada acesso gera \ev{support\_open} com \texttt{layer} $\in$ \{dica, ajuda, por\_que, caderno\}, \texttt{topic} (H1--H10 ou tela), \texttt{scene\_origin} e \texttt{step}; o fechamento gera \ev{support\_close} com \texttt{scene\_return}. Uso de suporte por camada: $U_{\ell} = N(\ev{support\_open}, \texttt{layer}=\ell) / N_{\mathrm{tarefas\ iniciadas}}$; $U = \sum_{\ell \neq \mathrm{caderno}} U_{\ell}$ (o caderno é relatado à parte, por ser só sintético); $U_{\mathrm{pre}}$ como na Fase~II. Abrir um tópico dentro de S6 não conta nova consulta; conta-se a abertura de S6.

\textit{Igualdade e congelamento.} O texto da ajuda é o mesmo para pessoas e agente, versionado por \textit{hash} junto da interface. A ajuda permitida é a mesma na calibração e na validação; na calibração, nenhum ajuste de conteúdo foi feito com base no desempenho do agente. \textit{Agente de ajuda não é UX-Agent}: o CasaFácil não tem assistente conversacional; se tivesse, ele seria parte do sistema e seguiria estas regras, enquanto o UXAgent está do outro lado e representa a persona.

% =====================================================================
\section{Extensões do esquema de eventos}\label{sec:eventos}

\begin{table}[!h]
\caption{Eventos e campos acrescentados ao esquema da Fase~II}
\label{tab:eventos}
\centering
\footnotesize
\begin{tabularx}{\textwidth}{@{}P{3.4cm}YP{3.2cm}@{}}
\toprule
\textbf{Evento ou campo} & \textbf{Conteúdo} & \textbf{Momento} \\
\midrule
\ev{dialog\_shown} & Tipo (AL, CE1, CE2), cena, texto exibido & Ao exibir o diálogo ou a mensagem \\
\ev{dialog\_choice} & Opção escolhida (excluir, cancelar, salvar, sair\_sem\_salvar) & Ao tocar na opção \\
\ev{support\_open}: \texttt{layer}, \texttt{topic}, \texttt{scene\_origin} & Camada, tópico e cena de origem & Ao abrir \\
\ev{support\_close}: \texttt{scene\_return} & Cena para onde voltou & Ao fechar \\
\ev{screen\_view}: \texttt{elements} & Lista dos signos visíveis na ordem visual (sem identificadores internos) & A cada exibição \\
\ev{method\_selected} (derivado) & Método usado por subobjetivo (historico|painel; interruptor|desfazer; pausar|editar) & Pós-processamento, pelo verificador \\
\ev{communicability\_tagged} (derivado) & Etiqueta MAC e eventos de origem & Pós-processamento, pelo verificador \\
\texttt{design\_version}, \texttt{help\_hash} & Versão do caderno e \textit{hash} do texto da ajuda & Em todo evento \\
\bottomrule
\end{tabularx}
\end{table}

% =====================================================================
\section{Registro de \textit{rationale}}\label{sec:rationale}

\begin{table}[!h]
\caption{Decisões de design (Q: questão; O: opções; C: critérios; D: decisão)~\cite{dix2004}}
\label{tab:rationale}
\centering
\scriptsize
\begin{tabularx}{\textwidth}{@{}P{0.5cm}P{3.0cm}P{3.4cm}YP{0.9cm}@{}}
\toprule
\textbf{Id} & \textbf{Questão} & \textbf{Opções} & \textbf{Critérios e decisão} & \textbf{Status} \\
\midrule
D1 & Como descrever a tela ao agente? & (a) árvore de acessibilidade; (b) lista com identificadores e estados; (c) só signos visíveis em texto literal & Equivalência de canal; sem vazamento. D: (c), regras da Seção~\ref{sec:signos} & \stPD \\
D2 & Usar botões desabilitados? & (a) sim, escrito em $o_t$; (b) sim, omitido da lista; (c) não usar & (a) vaza estado; (b) vaza que não funciona. D: (c), com CE onde seria natural desabilitar & \stPD \\
D3 & Como expressar ícones sem texto ao agente? & (a) função (``pausar''); (b) forma; (c) omitir & (a) vaza significado; (c) esconde o que a persona vê. D: (b), com M2 para a não interpretação & \stHO \\
D4 & Explicação na versão B & B1 sob demanda; B2 sempre exibida & Pergunta do estudo; observabilidade de $U$; densidade da tela; KLM: B2 2,73~s mais rápida. D: B1 & \stPD/\stHO \\
D5 & Proteger ``Excluir rotina'' & (a) sem proteção; (b) AL; (c) AL + RA; (d) PA (remover de S3) & Ação irreversível ao lado de ícones; receio de errar (L4)~\cite[p.~37]{tsuchiya2018}; (d) apagaria a oportunidade de E-ACT. D: (c) & \stEV/\stPD \\
D6 & Sair de S4 sem salvar & (a) descartar calado; (b) salvar automático; (c) CE1 & (b) mudaria a rotina sem intenção explícita; (a) gera E-OMI silencioso. D: (c), justificada por C2, não por V2 & \stPD \\
D7 & Confirmar sem resposta em S7 & (a) desabilitar; (b) pré-marcar ``Não sei''; (c) CE2 & (a) ver D2; (b) enviesa a resposta. D: (c) & \stPD \\
D8 & Camadas da ajuda & (a) só S6; (b) S6 + dica local; (c) ajuda adaptativa & Contextual e em camadas~\cite{barbosa2010}; (c) exigiria regras congeladas e mudaria $U$. D: (b), não adaptativa & \stPD \\
D9 & Conteúdo procedimental da ajuda & (a) genérico, do GOMS; (b) por tarefa & (b) entregaria a solução. D: (a) & \stPD \\
D10 & Limite do horário & (a) circular 23h$\rightarrow$0h; (b) preso em 00h--23h (PA) com nota & (a) gera erro sem sinal. D: (b), PA com PP & \stPD \\
D11 & Onde ficam Concluir e Desistir? & (a) só ao fim; (b) barra fixa & O protocolo permite encerrar a qualquer momento; E-OMI precisa ser possível. D: (b) & \stPD \\
D12 & Manter ícones sem texto e ``Smart Schedule''? & (a) trocar por texto em português; (b) manter & (a) seguiria a metamensagem, mas eliminaria L2, condição do estudo. D: (b), declarada como restrição; compensada pela ajuda & \stEV{} \cite[pp.~34--35]{tsuchiya2018} \\
D13 & Aviso de mudança em A e B & (a) só em B; (b) igual nas duas & Não confundir visibilidade da iniciativa com explicação~\cite{li2026}. D: (b) & \stPD \\
D14 & Estado do aparelho & (a) só interruptor; (b) interruptor + texto & Não depender de cor ou posição; texto literal em $o_t$ sem vazamento. D: (b) & \stPD \\
\bottomrule
\end{tabularx}

\end{table}

\textit{Verificação antes de gerar dados (a fazer na implementação).} (1)~Inspeção de vazamento: para cada tela, comparar a captura da tela com $o_t$ e conferir que nenhum termo de $o_t$ está ausente da tela. (2)~Lista de verificação da ajuda contra os signos de aplicação. (3)~Piloto de instrumentação com C1 e C2 apenas, para conferir que todo evento das Tabelas de eventos aparece e que o verificador decide o sucesso sem consultar o agente. (4)~Congelamento: \textit{hash} da interface, da ajuda, dos prompts e da configuração registrados antes de qualquer execução de V1 e V2.

% =====================================================================
\bibliographystyle{IEEEtran}
\bibliography{referencias}

\end{document}
